\cotefile{79}
\scanpage{235}
% Upper edge, fragment of the preceding strip
On trouve
\mspage{119 bis}

\unsure{graphes|graphe}, \inserted{avec} groupe d'homothéties $Z$ \unsure{qu'on se le donne|qu'on se donne|opérant librement}, \inserted{et satisfaisant la condition (41\unsure{*|''})}, \deleted{revient à celle} revient \deleted{à la} aux données \deleted{des choses} suivantes

\[
\left\{\begin{aligned}
&\text{a) un \deleted{gr.} }S^{\unsure{\sharp|\natural|\flat}},\ \text{une relation symétrique }R^\sharp\\
&\qquad\text{(i.e. un \uline{graphe} }\gamma^\sharp\text{) }\deleted{\unsure{\mathrm{sur}|\mathrm{entre}}}\ p_1,p_2:R^\sharp\rightrightarrows S^\sharp\\
&\text{b) }\deleted{\unsure{S}}\ \text{un }Z\text{-torseur relatif }S/S^\sharp\text{ sur }S^\sharp\\
&\text{c) un isomorphisme}\\
&\qquad p_1^*(S^{\deleted{\unsure{\sharp}}}/S^\sharp)\deleted{\unsure{\otimes}}\mathbin{\wedge}p_2^*(S/S^\sharp)\simeq\underset{\text{torseur trivial}}{\mathbb{1}_{Z_{S^\sharp}}}
\end{aligned}\right.
\tag*{(48.7)}
\]

% slanted
\margin{C'est \deleted{\unsure{essentiellement}} la relation symétrique sur $S$ associée à la structure de graphe}

Enfin, dans cette optique, la donnée d'un multigraphe $\pi$, \deleted{\unsure{avec|$\pi$}} avec groupe d'homothéties, revient aux données suivantes

\[
\left\{\begin{aligned}
&\text{a) Un graphe }\gamma^\sharp=(S^\sharp,\underbrace{R^\sharp\subset S^\sharp\times S^\sharp}_{\text{rel. sym.}})\\
&\text{b) Un }Z\text{-torseur relatif }S^{\deleted{\unsure{\sharp}}}/S^\sharp\\
&\text{c) Un }Z\text{-torseur absolu }\pi\\
&\text{d) un isomorphisme}\\
&\qquad p_1^*(S/S^\sharp)\wedge p_2^*(S/S^\sharp)\simeq\pi
\end{aligned}\right.
\tag*{(48.8)}
\]

Ainsi : la description (49.7), dont $i\in\pi$ définit alors une structure de graphes à homothéties sur $S$, comme ens de sommets, puisqu'il définit un iso $p_1^*(S/S^{\unsure{\sharp|\flat}})\wedge p_2^*(S/S^\sharp)\simeq\pi\xrightarrow[\text{via }i]{\sim}\mathbb{1}$.

\uline{Ex.} Si $M$ module libre de $\mathrm{rg}\,2$ sur $\Lambda$, $S=M^*/{\pm1}$, $\check{S}^{*}\simeq\check{M}^{*}/{\pm1}$ si $Z\deleted{\unsure{\mathrm{anneau}}}=\Lambda^*/{\pm1}$, \deleted{\unsure{on}} $\pi=(\det M)^{*}/{\pm1}$, on a $S^\sharp=M^*/\Lambda^*\subset\mathbb{P}(M)(\Lambda)$ (égalité si $\mathrm{Pic}(\Lambda)\simeq\unsure{0|0'}$), $R^\sharp\deleted{\unsure{=}}\subset S^\sharp\times S^\sharp$ est formé des \deleted{classes de couples} \deleted{\unsure{$(x,y)$}} $(x^\sharp,y^\sharp)$, avec $x,y\in M^*$, $x\wedge y\in(\det)^{*}$ --- i.e. les couples de sections du fibré projectif, distinctes en tt pt de la fibre. Le $S$-torseur relatif, sur \deleted{$\mathbb{P}(M)$} $S^\sharp$ correspond au fibré des \unsure{Serre-Hopf|semi-droites} modulo $\pm1$. L'iso. \deleted{\unsure{est}} d) est celui déduit de $\det M^*\simeq D^{\unsure{\wedge|1|*}}\mathbin{\wedge}_{\Lambda^*}D'^{\unsure{\wedge|1|*}}$ en passant à $\Lambda^*/{\pm1}$

% slanted
\margin{ou encore les couples de \textquotedblleft{}droites\textquotedblright{} $D,D'$ dans $M$ telles que $M=D\oplus D'$ --- d'où $\deleted{\unsure{\det}}\det M\simeq D\otimes D'$ donc $\det M^*\simeq D^\wedge\mathbin{\wedge}_{\Lambda^*}D'^\wedge$}

% archive sheet number: 234
\scanpage{236}
\mspage{$120'$}

2. Condition d'isomorphie pour une donnée $\pi$.

\deleted{On introduit}

\[
\deleted{(45)\qquad G=\{g\in A=\mathrm{Aut}(\Gamma)\mid \substack{g\text{ centralise }Z\ \text{i.e. }g\in\mathcal G\\g_{\unsure{\mathcal G}}(\pi)=}\,}
\]

On s'intéresse aux automorphismes d'une structure \deleted{$\underset{\substack{\shortparallel\\(S,\check S,R)}}{\Gamma},Z\subset\unsure{G}_{S}\ \unsure{Q}_{\check S}$}

\[
M=\bigl(\underset{\substack{\shortparallel\\(S,\check S,R}}{\Gamma},\quad Z\subset\mathcal A=\mathrm{Aut}(\Gamma),\quad\pi\subset\Pi_Z\subset\mathrm{Bij}(S,\check S)\bigr),
\tag*{(49)}
\]

% slanted
\margin{\textquotedblleft{}\unsure{prégraphisme|prégraphe} \deleted{\unsure{groupe}} multijaugé\textquotedblright}

ce sont donc les $g\in\mathcal A$ qui \uline{normalisent} $Z$, et qui satisfont

\[
\underbrace{g(\pi)}_{\text{tr. de str.}}\left(\overset{\mathrm{def}}{=}\underbrace{\{g_{\check S}\circ\varphi\circ g_S^{-1}\mid\varphi\in\pi\}}\right)=\pi
\tag*{(50)}
\]

ce qui équivaut à la condition, \deleted{\unsure{que}} \inserted{en termes d'\uline{un} $\varphi_0\in\pi$}

\[
g(\varphi_0)\left(\overset{\mathrm{def}}{=}g_{\check S}\varphi_0g_{\check S}\varphi^{-1}\right)\in\pi\quad.
\tag*{(51)}
\]

On a donc des hom de groupes

\diagram{\begin{tikzcd}[ampersand replacement=\&,column sep=large]
{(52)}\&\mathrm{Aut}(M)\arrow[r]\&\mathrm{Aff}_{Z}(\pi)\arrow[r]\&\mathrm{Aut}(Z),\\
\&\mathcal A\arrow[u,hook]\&\&
\end{tikzcd}}

et je m'intéresse aux autom. $g$ de $M$ dont l'image dans $\mathrm{Aut}(Z)$ soit égale à $1$, i.e. \inserted{qui \uline{centralisent} $Z$} i.e. qui opèrent sur $\pi$ par translation, \deleted{On} donc on s'intéresse au groupe formé de ces $g$ :

\[
G=\left\{g\in\mathcal G\overset{\mathrm{def}}{=}\mathrm{Centr}_{\mathcal A}(Z)\ \middle|\ \substack{g(\pi)=\pi\quad\mathrm{i.e.}\\g_{\unsure{S|\check S}}\varphi_0g_S^{-1}=\varphi_0}\right\}.
\tag*{(53)}
\]

On trouve donc un diagramme canonique.

% archive sheet number: 235

\mspage{119 bis}
\unsure{graphes|graphe}, \inserted{avec} groupe d'homothéties $Z$ \unsure{qu'on se le donne|qu'on se donne|opérant librement}, \inserted{et satisfaisant la condition (41\unsure{*|''})}, \deleted{revient à celle}
\scanpage{237}
\mspage{121}

\diagram{\begin{tikzcd}[ampersand replacement=\&]
{(54)}\&1\arrow[r]\&G_0\arrow[r]\&G\arrow[r,"\delta"]\&Z\\
\&\&\&Z\arrow[u,hook]\arrow[ur,"{\lambda\mapsto\lambda^2}"']\&{,}
\end{tikzcd}}

où $\delta$ (jouant le rôle d'un \textquotedblleft{}déterminant\textquotedblright) exprime l'opération de $G$ sur $\pi$ par translations, et où on a posé

\[
\begin{aligned}
G_0=\mathrm{Ker}\,\delta&=\{g\in\deleted{\mathcal G}\mathcal A\mid g(\varphi)=\varphi\quad\forall\varphi\in\pi\}\\
&=\left\{g\in\mathcal A\ \middle|\ \substack{\deleted{\text{i.e.}}\ g\in\mathcal G\ \text{i.e. }g\text{ centralise }Z\\g(\varphi_0)=\varphi_0}\right\}.
\end{aligned}
\tag*{(55)}
\]

\margin{\uline{NB} Si $R$ séparée en $\check S$, on a $(51')\ \left\{\begin{aligned}G_0&\xrightarrow{\sim}\text{groupe des automorphismes}\\[-2pt]&\text{du graphe }\gamma_{\varphi_0}=(S,A_{\varphi_0})\\ \lambda&\longmapsto\lambda_S\quad\text{qui commutent à }Z_S.\end{aligned}\right.$}

Le fait que $Z\subset G$, et la commutativité exprimée dans (54)

\[
\delta(\lambda)=\lambda^2\qquad\forall\lambda\in Z
\tag*{(56)}
\]

provient de la relation

\[
\lambda(\varphi)=\lambda_{\check S}\circ\underbrace{\varphi\circ\lambda_S^{-1}}_{\lambda_{\check S}\varphi}=\lambda_{\check S}^{2}\circ\varphi=T_{\lambda^2}(\varphi).
\]

La définition de $\delta(g)$ pour $g\in G$ s'explicite par

\[
g(\varphi)=T_{\delta(g)}\varphi\ \bigl(=\delta(g)_{\check S}\circ\varphi=\varphi\circ\delta(g)_S\bigr)\qquad\substack{\forall\varphi\in\pi\\g\in G}
\tag*{(57)}
\]

qui caractérise \deleted{\unsure{$\forall$}} $\delta(g)$, pour \uline{un} $\varphi=\varphi_0$ fixé déjà.

On a les Conditions équivalentes sur $M=(\Gamma,Z,\pi)$

\[
\begin{aligned}
&\text{a) L'hom }\delta\text{ est surjectif}\\
&\text{b) }\deleted{\unsure{\mathcal G}}G\text{ opère transitivement sur }\pi\\
&\text{c) }\forall\lambda\in Z,\ \deleted{\unsure{\lambda\varphi_0}}\ \text{la structure de}\\
&\quad\text{graphes définie par } A_\lambda\subset\underset{\{1,2\}}{\mathfrak P}(S)\\
&\quad A_\lambda=\{\{s,\lambda t\}\mid\{s,t\}\in S,\ \text{t.q. }\{s,\deleted{\unsure{\lambda}}t\}\inA_0\}\\
&\quad\text{est isomorphe à }\unsure{\gamma}(S,A_0),\ \text{par une}\\
&\quad\deleted{\unsure{f}}\ f\in\unsure{\mathfrak S|\mathcal G}_S\text{ qui centralise }Z_S\\
&\text{d) }\inserted{\forall\lambda\in Z}\ \exists f_\lambda\in\mathfrak S_S,\ \text{centralisant }Z_S,\ \text{telle que}\\
&\quad\forall s,t\in S,\ \text{on a }\{s,t\}\in A_0\Longleftrightarrow\{f_\lambda^{-1}s,f_\lambda^{-1}\lambda t\}\in A_0
\end{aligned}
\tag*{(58)}
\]

% slanted
\margin{ici on suppose choisi un $\varphi_0\in\pi$, d'où une structure de graphes --- et on suppose $R$ \deleted{\unsure{non}} séparée en $\check S$}

On dit que la condition d'homogénéité est satisfaite

% archive sheet number: 236
\scanpage{238}
si l'hom $\delta$ est surjectif, ie les conditions (\unsure{59|58}) satisfaites.

\deleted{On}

% archive sheet number: 237
\scanpage{239}
\mspage{122}

3. Conditions d'homogénéité.

\deleted{Nous nous intéressons maintenant au cas où} Revenons au cas de \deleted{$G_0$ a la propriété d'homogénéité} \deleted{\unsure{$\Gamma$}}

% slanted
\margin{Cas d'un}

\[
M=(\Gamma,\quad Z,\quad\pi),
\]

je dis qu'on a équivalence entre les propriétés

\[
\left\{\begin{aligned}
&\text{a) }G_0\text{ est transitif sur }R\\
&\text{b) }G_0\text{ est transitif sur }R_{\varphi_0}\ \text{i.e. sur l'ens }\overrightarrow{A_{\varphi_0}}\\
&\qquad\text{des arêtes orientées du graphe }(S,A_{\varphi_0}),\ \text{pour}\\
&\qquad\text{\uline{un} }\varphi_0\in\pi\\
&\text{b$'$) Comme b), pour tout }\varphi\in\pi
\end{aligned}\right.
\tag*{(59)}
\]

% slanted
\margin{\uline{NB} a) s'énonce sans référence à un $\pi$ contrairement à b) b$'$) --- c'est une propriété d'un \deleted{\unsure{$\pi$ mais}} groupe d'autom $G_0$ d'un prégraphe, en \deleted{fait} \inserted{\unsure{présence|défini}} d'un $Z$. (Mais on suppose ici $G_0$ en termes d'un $Z,\pi$).}

L'équivalence des a) et b) \inserted{ou a) et b$'$)} résulte des \deleted{\unsure{ceci} fait}

\[
R\simeq R_{\varphi_0}\qquad(G_0\text{-isom})\quad\text{via }\mathrm{id}\times\varphi_0\,\unsure{.}
\]

\deleted{cette} \deleted{\unsure{du b) et}} Nous dirons que la situation \deleted{du} multigraphe $M$ est \uline{homogène}, s'il satisfait (59) \inserted{et si \uuline{de plus}} \deleted{les conditions indiquées, de préférence si $R\ne\varnothing$,} \deleted{que} $G_0$ est transitif sur $S$ et sur $\check S$ \unsure{.|;} \deleted{et par suite} \deleted{suite que si} \inserted{alors} un $\lambda\in Z$ \deleted{\unsure{opère}} a un pt fixe dans $S$ (resp $\check S$), il opère trivialement dans $S$ (resp. ds $\check S$), donc il est trivial si $\Gamma$ est séparé en $\check S$ (resp. en $S$). Donc

(60) Si $\Gamma$ séparé en $\check S$ (resp $S^{\unsure{*}}$), alors sous les conditions d'homogénéité (55), $Z$ opère librement sur $S$ (resp. $\check S$).

On a aussi

(61) (55) implique (41), \deleted{\unsure{avec les}} \inserted{et ses} formes équivalentes $(41')\ (41'')\ (44)\ (48)$ --- donnant donc lieu aux interprétations (48.1) à (48.8.)

\deleted{(avec 59\unsure{.|a}) Ce qui s'écrit encore (sans référence à $\pi$, en termes seulement de $G_0,Z\subset\mathcal A$ deux sous-groupes de $\mathcal A=\mathrm{Aut}\,\Gamma$ commutant l'un à l'autre, sous la forme 41}
\[
\deleted{(\unsure{*})\qquad (s_0,\check s_0)\in R\text{ et }(s_0,\lambda\check s_0)\in R,\quad\lambda\in Z\quad\Longrightarrow\lambda=1}
\]
\deleted{en effet il existe $g\in G_0$ tel que $g(s_0,\check s_0)=(s_0,\lambda\check s_0)$,}

% archive sheet number: 238
\scanpage{240}
\deleted{i.e.}
\[
\deleted{gs_0=s_0,\quad g\check s_0=\lambda\check s_0\quad\text{i.e. }\unsure{g}\check s_{0}=g^{-1}\lambda\check s_0=\lambda g^{-1}\check s_0}
\]

en effet, comme $G_0$ opère transitivement sur $A_0$ et $A_\lambda$ (il suffit \inserted{donc} pour cet argument que l'on ait cette transitivité \textquotedblleft{}faible\textquotedblright, au lieu de la transitivité sur $\overrightarrow{A_0}$ \deleted{\unsure{fibre sur}} i.e. sur $R$), si on \deleted{\unsure{voit}} $A_0\cap A_\lambda=\varnothing$, on en conclut $A_0=A_\lambda$ d'où $\lambda=\mathrm{id}$, cqfd.

Nous voulons maintenant expliciter \deleted{\unsure{les diverses}} \inserted{la} \deleted{\unsure{situations|donnée}} \inserted{donnée} d'un \deleted{\unsure{de}} multigraphe à homothéties homogène (i.e. satisfaisant (59)), en termes de données de théorie des groupes --- plus particulièrement dans le cas où on suppose de plus satisfaite la condition d'isomorphie (58).

Pour ceci, \deleted{dans le \unsure{premier}} je vais prendre les choses par un bout légèrement différent, en \deleted{\unsure{me}} \uline{partant} d'un cas \textquotedblleft{}épinglé\textquotedblright, \deleted{graphe} savoir

\[
\left\{\begin{aligned}
&\text{a) graphe }\deleted{\unsure{\gamma}}\ \gamma_0=(S,A_0),\text{ muni d'\deleted{une arête}}\\
&\text{b) groupe }\deleted{\unsure{G_0}}H\text{ d'autom de }(S,A_0),\text{ transitif sur }\overrightarrow{A_0}\\
&\quad\inserted{[\text{\uline{et} sur }S,\ \text{i.e. }\forall s\in S,\ \mathrm{Et}\ A_0(s)\ne\varnothing]}\\
&\text{c) et d'une \textquotedblleft{}origine\textquotedblright{} }(s_0,t_0),\text{ avec donc }\{s_0,t_0\}\in A_0\\
&\quad\text{et enfin}\\
&\text{d) d'un sous-groupe } Z\subset\unsure{\mathfrak S|\mathcal G}_S\text{ commutatif,}\\
&\quad\text{formé d'homothéties pour le graphe }\gamma_0\\
&\quad\text{commutant à l'action de }\deleted{\unsure{G_0}}H,
\end{aligned}\right.
\tag*{(62)}
\]

sans nous préoccuper pour le moment \unsure{si|où} \deleted{\unsure{$G_0$}}\inserted{$H$} est \inserted{le groupe $G_0$}, formé de \uline{tous} les autom. \deleted{qui commutent} \inserted{de la structure} qui commutent à $Z$ --- et en fait je serai amené à travailler avec la donnée d'un sous-groupe \deleted{\unsure{du}} plus petit, $G_0^{+}$. \deleted{Autant alors partir tout de suite d'un tel sous-groupe, je vais le noter $H$, pour identifier la donnée correspondante. On voit que}

% archive sheet number: 239
\scanpage{241}
\mspage{123}

\deleted{la subtilité que \inserted{voudrai} \unsure{j'introduirai|j'introduirais} dans mon groupe $G$}

Je considère cet $H$ comme une structure supplémentaire, en plus du \deleted{$Z$} groupe d'homothéties $Z$, et définirai $G$ en conséquence par

\[
\begin{aligned}
G={}&\deleted{\text{groupe des}}\ \bigl\{f\in\unsure{\mathfrak S|\mathcal G}_S\mid f\text{ normalise }G,\text{ centralise}\\
& Z,\ \text{et }\exists\lambda\in Z\text{ tel}\\
&\text{que } f(A_0)=A_\lambda\quad\text{i.e.}\\
&\forall s,t\in S,\quad\{s,t\}\in A_0\Longleftrightarrow\{f^{-1}s,\unsure{f\lambda^{-1}t|f^{-1}\lambda t}\}\in A_0
\end{aligned}
\tag*{(63)}
\]

de sorte qu'on aura \deleted{\unsure{encore}} un diagramme

\diagram{\begin{tikzcd}[ampersand replacement=\&]
\&\&H\arrow[d,hook]\arrow[dr,dashed,"\text{\uline{distingué}}"]\&\&\\
{(64)}\&1\arrow[r]\&\underset{\substack{\shortparallel\\\mathrm{Ker}\,\delta}}{G_0}\arrow[r]\&G\arrow[r,"\delta",no head]\&Z\\
\&\&\&Z\arrow[u,hook]\arrow[ur,"{\lambda\mapsto\lambda^2}"']\&
\end{tikzcd}}

\margin{$\nearrow$ Pour écrire ici $Z$, et non un quotient de $Z$, on admet que $Z$ opère \uline{fidèlement} sur $\pi$, ce qui est automatique si tt sommet de $\gamma_0$ est déterminé par son étoile}

où l'inclusion composée $H\longrightarrow G_0\unsure{-|\to}G$ fait de $H$ un sous-groupe \uline{distingué} de $G$. Pour l'instant on ne suppose pas $\delta$ surjectif. On fera attention que dans cette présentation, même si on suppose $\delta$ surjectif, les groupes $G_0$, \inserted{$G$} introduits ici \deleted{\unsure{ne sont pas}} \inserted{sont à priori} plus petits que ceux du n$^\circ$ précédent, qui eux ne font pas intervenir la donnée particulière \unsure{du hom|du choix|d'un} $H$. En fait, ceux \deleted{derniers \unsure{deux}} \inserted{envisagés ici sont} les \uline{normalisateurs} de $H$ dans les groupes $G_0^{\unsure{!}},G$ envisagés au n$^\circ$ précédent. L'on retrouvera \deleted{la situation} ces derniers, bien sûr, en prenant $H=G_0$.

Notre objet principal est de décrire de façon aussi précise que possible le diagramme (64),

% archive sheet number: 240
\scanpage{242}
et toute la situation de départ (62), en termes de \deleted{données} données de théorie des groupes plus explicites ou \textquotedblleft{}élémentaires\textquotedblright.

Quand on part de l'ensemble $S$, avec les deux sous-groupes

\[
\deleted{\unsure{G}}H\deleted{\unsure{0|Z|\sim}},Z\subset\unsure{\mathfrak S|\mathcal G}_S
\tag*{(65)}
\]

qui commutent l'un à l'autre, $Z$ commutatif, alors \deleted{la} \inserted{une} structure de \inserted{multi}graphe \inserted{$(S,\deleted{\unsure{A_0|A=}}(A_i)_{i\in\pi})$} sur $S$, admettant $H$ comme groupe d'automorphismes transitif \deleted{\unsure{sur les objets}} sur \inserted{les} $\overrightarrow{A_i}$ \inserted{$(i\in\pi)$}, et $Z$ comme groupe d'homothéties, est connue quand on connaît \inserted{($\pi$, et)} $A_{\bullet}(=\bigcup_{i\in\pi}A_i)\subset\mathfrak P_{\{1,2\}}(S)$,

\[
A\subset\underbrace{\mathfrak P_{\{1,2\}}}_{\substack{\text{parties à }1\\\text{ou }2\text{ élém.}}}(S),
\tag*{(66)}
\]

\deleted{qui doit être nécessairement un espace homogène} Car on récupère alors les $A_i$ $(i\in\pi)$ comme \deleted{\unsure{sous les groupes|sous le groupe}} les orbites de $H$ dans $A$, et on aura

\[
\deleted{(67)\qquad\widetilde H=H\mathbin{.} Z\subset\mathfrak S_S\quad -}
\]

\[
\pi\simeq A/H\quad .
\tag*{(67.)}
\]

\deleted{car on retrouve alors} \deleted{Pour \deleted{Explicitons} les conditions sur les données \inserted{Nous explicitons} (65) (66) pour qu'elles correspondent bien à un multigraphe \deleted{homogène} à groupe \deleted{d'automorphe} \inserted{d'homothéties} $Z$ et à groupe d'automorphismes homogène $A$, on trouve \deleted{que} \inserted{d'abord}}

\deleted{(67) $A$ stable par \unsure{$\sigma$} i.e. $H$ est un groupe d'autom pour $(S,A)$}

\[
\deleted{(68)\qquad\forall(s,t)\in S\times S,\ \inserted{\lambda\in Z,}\ \{s,t\}\in A\Longrightarrow\{\lambda s,t\}\in A}
\]

\deleted{--- condition qui est bien sûr valable pour $H$}

% archive sheet number: 241
\scanpage{243}
\mspage{124}

Il serait facile d'expliciter les conditions sur les données (65) (66) pour qu'elles correspondent à un multigraphe à homothéties et gr. d'automs homogène $H$ --- je préfère plutôt \unsure{d\textquoteright{}écrire|écrire} les conditions pour que les données

\[
H,Z\subset\unsure{\mathfrak S|\mathcal G}_S\quad\text{et}\quad (s_0,t_0)\in S\times S
\tag*{(68)}
\]

correspondent à un multigraphe à homothéties et groupe d'autom \deleted{\unsure{pour}}\inserted{hom} --- en notant que la donnée de $(s_0,t_0)\in A$ \deleted{permet}, en plus de $H,Z$, permet de récupérer $A$, comme

\[
A=\bigl\{\{u\unsure{\lambda|(}s_0,u\lambda t_0\}\mid\lambda\in Z,\ u\in\deleted{\unsure{G|O}}H\bigr\}.
\tag*{(69)}
\]

On trouve le système de conditions

\[
\left\{\begin{aligned}
&\text{a) }H\text{ transitif sur }S,\ H\text{ et }Z\text{ commutent.}\\
&\text{b) }\forall\lambda\in Z,\ \exists u\in H\quad\text{tel que}\\
&\qquad u(s_0,t_0)=(\lambda^{-1}s_0,\lambda t_0).
\end{aligned}\right.
\tag*{(70)}
\]

On récupère alors \deleted{$A_0$}

\[
A_0=\bigl\{\{u\deleted{\unsure{s|\lambda}}s_0,ut_0\}\mid u\in H\bigr\},
\tag*{(\unsure{78|71})}
\]

et par b) $Z$ sera un groupe d'homothéties pour $(S,A_0)$, ce qui engendre un multigraphe à homothéties $(S,A)$ :

On n'exclut pas le cas où on aurait $s_0=t_0$, auquel cas $A_0$ est formé de toutes les parties ponctuelles de $S$, et la condition b) signifie que $\lambda\in Z\Longrightarrow\lambda^2=1$. Si $\lambda\in Z$, $\lambda\ne1$, alors $A_\lambda$ est formé des orbites sous $\lambda$, qui sont des parties à deux éléments. \deleted{L'étoile}

% archive sheet number: 242
\scanpage{244}
Si on ne fait pas d'hypothèse supplémentaire, il n'est pas sûr que l'opération de $Z$ sur $\pi$ soit fidèle, i.e. qu'on ait

\[
\left\{\begin{aligned}
&\forall\lambda\in Z\ \deleted{\text{tel qu'il existe}}\ u\in H\ \deleted{\text{avec}}\\
&(s_0,\lambda t_0)=u(s_0,t_0)\qquad\Longrightarrow\lambda=1\quad.\\
&\text{i.e. }\quad us_0=s_0,\ u(t_0)=\lambda t_0
\end{aligned}\right.
\tag*{(72)}
\]

Pour ceci, il suffit que tout sommet soit déterminé par son étoile --- condition qui ne s'exprime pas, il me semble, de façon particulièrement jolie en termes des données $\unsure{S|Z};H\subset\unsure{\mathfrak S|\mathcal G}_S,\ (s_0,t_0)\in S\times S$.

Je vais, pour simplifier, supposer qu'on a (72), --- sinon dans (64), l'hom $\unsure{\delta_1|\delta,}$ va de \deleted{\unsure{$H$}} $G$, non dans $Z$, mais dans le groupe quotient $Z/\unsure{\mathfrak z|\zeta}_Z$, où $\mathfrak z_Z$ est le noyau de l'opération de $Z$ sur $\pi$, i.e le sous-groupe formé des $\lambda\in Z$ pour lesquels $\exists u\in H$ avec $u(s_0,t_0)=(s_0,\lambda t_0)$.

Posons

\[
\left\{\begin{aligned}
U_0&=\deleted{\unsure{G|O}}H_{s_0}=\{u\in H\mid us_0=s_0\}\\
U_1&=H_{t_0}=\{u\in H\mid ut_0=\unsure{t|t_0}\}\\
\mathfrak z_H&=U_0\cap U_1=H_{(s_0,t_0)}=\{u\in H\mid us_0=s_0,ut_0=t_0\,,
\end{aligned}\right.
\tag*{(73)}
\]

et soit

\[
\sigma\in H\qquad\sigma(s_0,t_0)=(t_0,s_0)\quad\text{i.e.}\quad\sigma s_0=t_0,\sigma t_0=s_0.
\tag*{(74)}
\]

On a donc

\[
\left\{\begin{aligned}
U_1&=\sigma(U_0),\quad U_0=\sigma(\deleted{\unsure{H}}U_1)\\
\sigma^2&\in\mathfrak z_H\overset{\mathrm{def}}{=}U_0\cap U_1\ \deleted{\unsure{\sim}}\quad.
\end{aligned}\right.
\tag*{(75)}
\]

% archive sheet number: 243
\scanpage{245}
\mspage{125}

On voit tout de suite

(76) La donnée du graphe à opérateurs homogène $(S,A_0\underset{\deleted{\unsure{\gg|=}}}{\subset}\mathfrak P_{\{1,2\}}(S), H\subset\unsure{\mathfrak S|\mathcal G}_S)$, muni de l'arête orientée $(s_0,t_0)$, équivaut à la donnée du groupe $H$, des deux sous-groupes $U_0,U_1$ de $H$, et de $\sigma\in H$ satisfaisant les relations (75), et de plus

\[
\bigcap_{u\in H}\mathrm{int}(u)U_0=\{1\}
\tag*{(77)}
\]

qui exprime que l'opération de $H$ sur l'espace homogène $H/U_0$ est fidèle.

On récupère la situation géométrique, \unsure{et|de|ou} la situation algébrique, par

\[
\begin{aligned}
&S\simeq H/U_0,\quad s_0=U_0,\quad t_0=\sigma\mathbin{.}H_0\quad\text{(translaté à gauche)}\\
&\overrightarrow{A_0}\simeq H/\unsure{\mathfrak z|\zeta}
\end{aligned}
\tag*{(78)}
\]

\deleted{Dans} L'énoncé (76) n'est correct que si, ou bien on sous-entend que parmi les données géométriques il y a aussi le \uline{choix} de $\sigma$ satisfaisant (74), soit que, dans la traduction algébrique, $\sigma$ \deleted{n'est donné que} \inserted{(qui normalise $\mathfrak z_H$, nécessairement)} n'est donné que \textquotedblleft{}modulo $\mathfrak z_H$\textquotedblright{} (i.e. que ce qui est donné est la classe \inserted{$\sigma\mathfrak z_H=\mathfrak z_H\sigma$} \deleted{de} $\sigma$ dans $\mathrm{Norm}_H(\mathfrak z_H)/\mathfrak z_H$). Cette ambiguïté \unsure{disparait|disparaît} quand on suppose que $H$ est \uline{simplement} transitif sur l'ens \inserted{$\overrightarrow{A_0}$} des arêtes orientées ; \deleted{ce que}

\[
\begin{aligned}
&H\text{ simplement transitif sur l'ens }\overrightarrow{A_0}\\
&\Longleftrightarrow\mathfrak z_H(\overset{\mathrm{def}}{=}U_0\cap U_1)=\{1\},\quad\deleted{\text{ce q}}
\end{aligned}
\tag*{(79)}
\]

Condition qui implique l'unicité de $\sigma$, et

\[
\sigma^2=1 .
\tag*{(80)}
\]

% archive sheet number: 244
\scanpage{246}
. Nous allons expliciter maintenant la donnée de $Z$, en introduisant le sous groupe de $H$

\[
\deleted{\unsure{D}}\ T_H\overset{\mathrm{def}}{=}\{\,\deleted{\unsure{t}}u\in H\mid\exists\lambda\in Z,\ \deleted{\unsure{t}}us_0=\lambda s_0,\ ut_0=\lambda^{-1}t_0\}
\tag*{(\unsure{80|88})}
\]

C'est bien un sous-groupe, \deleted{puisque si $u$ et} si $u$ correspond à $\lambda$, $v$ correspond à $\mu$, alors $vu$ correspond à $\mu\lambda$. \deleted{D'ailleurs le $\lambda$ associé} à $u$ est unique, car \deleted{\unsure{on est ramené}} \inserted{$Z$ opère \uline{librement} sur $S$} \deleted{voir que} \deleted{mieux}

\[
\deleted{(\lambda s_0,\lambda^{-1}t_0)=(s_0,t_0)\Longrightarrow\lambda=1\quad(\lambda\in Z)}
\]

\deleted{ce qui résulte aussitôt de (72).} On a donc un hom de groupes canonique $T_H\longrightarrow Z$, dont le noyau est $\unsure{\overline{\mathfrak z}|\overline\zeta|\mathfrak z_H}$,

\diagram{\begin{tikzcd}[ampersand replacement=\&]
{(81)}\&1\arrow[r,no head]\&\mathfrak z_H\arrow[r,arrows={-|}]\&T_H\arrow[r,"{\unsure{\kappa|\pi|k}}"]\&Z\arrow[r]\&1
\end{tikzcd}}

\margin{Attention, cet hom \uline{n'est pas} l'hom $\delta\mid Z$}

qui dans le cas simplement transitif s'écrit

\[
T_H\xrightarrow{\sim}Z\quad .
\tag*{$(81')$}
\]

Notons

\[
T_H\text{ normalise }U_0,\ \inserted{\text{et }U_1,}\quad T_H\cap U_0=T_H\cap U_1=\mathfrak z_H
\tag*{(82)}
\]

\[
\begin{aligned}
&\deleted{t\in T_H\Longrightarrow\sigma(t)\ \unsure{=}\ t^{-1}\quad\text{i.e. }\sigma\text{ normalise }T_H}\\
&\deleted{\text{et induit sur }T_H/\unsure{\mathfrak z_H}\ \unsure{\simeq}\ \text{l'aut}}\\
&\deleted{\text{du passage à l'inverse}}
\end{aligned}
\]

En effet, si $\deleted{\unsure{u}}t\in T_H$, $u\in U_0$, on a

\[
tut^{-1}(s_0)=t\underbrace{u\,\lambda^{-1}(s_0)}_{\lambda^{-1}\underbrace{u(s_0)}_{s_0}}=t(\lambda^{-1}s_0)=\lambda^{-1}ts_0=\lambda^{-1}\lambda s_0=s_0
\]

d'où $tut^{-1}\in U_0$, \deleted{\unsure{de même on}} \inserted{par symétrie on} voit \deleted{\unsure{tout}} que $T_H$ normalise $U_1$. Si $t\in T_H\cap U_0$, on a

\[
t(s_0,t_0)=(\lambda s_0,\lambda^{-1}t_0)=(s_0,\lambda^{-1}t_0)\qquad\text{puisque }ts_0=s_0
\]

% archive sheet number: 245
\scanpage{247}
\mspage{126}

d'où $\lambda=1$ par (72), d'où $t(s_0,t_0)=(s_0,t_0)$ i.e. $t\in\unsure{\mathfrak z|\zeta}_H$, donc $T_H\cap U_0\subset\mathfrak z_H$, l'\deleted{\unsure{impl}}inclusion inverse étant triviale. \deleted{Notons} Posons alors

\[
\left\{\begin{aligned}
B_{0H}&=T_H\mathbin{.}U_0\\
B_{1H}&=T_H\mathbin{.}U_1
\end{aligned}\right|\quad(\text{produits }1/2\text{-directs si }\mathfrak z_H=\{1\}),
\tag*{(83)}
\]

je dis qu'on a

\[
\left\{\begin{aligned}
B_{0H}&=\{u\in H\mid us_0\in Z\mathbin{.}s_0\}=H_{D_0}\\
B_{1H}&=\{u\in H\mid us_1\in Z\mathbin{.}\deleted{\unsure{s}}t_0\}=H_{D_1},
\end{aligned}\right.
\tag*{(84)}
\]

où $D_0,D_1$ sont les \textquotedblleft{}\uline{rayons}\textquotedblright{} de $s_0,t_0$ i.e leurs orbites sous $Z$. La vérification est immédiate du fait que l'hom $k:T_H\unsure{-|\to}Z$ est épi. Notons aussi

\[
B_{0H}\cap B_{1H}=T_H,\qquad\text{d'où \inserted{aussi} }B_{0H}\cap U_1=U_0\cap B_{1H}=\mathfrak z_H,
\tag*{(85)}
\]

en effet si $\unsure{k|t}\in B_{0H}\cap B_{1H}$, i.e. si $t\in H$ tel que l'on ait

\[
t(s_0,t_0)=(\lambda s_0,\mu s_0)\qquad\lambda,\mu\in Z,
\]

comme $(\lambda s_0,\mu s_0)$ est une $\lambda\mu$-arête, et $t\in H$\unsure{,} donc c'est une $A_0$-arête, on a $\lambda\mu=1$, i.e. $\mu=\lambda^{-1}$, d'où $t\in T$. \deleted{\unsure{Not}}

Notons maintenant que les groupes \deleted{\unsure{$T_0$}} $U_0,U_1,T_H,B_{0H},B_{1H}$ dépendent du choix de $(s_0,t_0)\in\overrightarrow{A_0}$ Si on échange $s_0$ et $t_0$, alors $U_0$ et $U_1$ s'interchangent, ainsi que $B_{0H}$ et $B_{1H}$, tandis que $T_H$ reste inchangé, il ne dépend que de \uline{l'arête} $\{s_0,t_0\}$ --- mais $k$ est remplacé par $k'=\deleted{\unsure{t}}t\mapsto k(t)^{-1}$.

% archive sheet number: 246
\scanpage{248}
En termes de $\sigma\in H$, \unsure{cel|cela} s'exprime par

\[
\left\{\begin{aligned}
&\sigma\text{ normalise }T_H,\text{ et }\deleted{\mathrm{int}(\sigma)\mid T_H/\unsure{\mathfrak z|\zeta}}\ \forall t\in T_H,\text{ on a}\\
&\qquad\sigma(t)\equiv t^{-1}\quad(\mathfrak z_H)\\
&\text{donc }\sigma(t)=t^{-1}\quad\text{si }\mathfrak z_H=\{1\}\quad.
\end{aligned}\right.
\tag*{(86)}
\]

\[
B_{1H}=\sigma(B_{0H}),\quad B_{0H}=\sigma(B_{1H})\quad.
\tag*{(87)}
\]

Notons maintenant

(88) La donnée, \deleted{de} pour un graphe homogène défini par $(H,U_0\subset H,\sigma\in H)$ (satisfaisant $\sigma^2\in U_0\cap\sigma(U_0)=\mathfrak z$) d'un groupe d'homothéties $Z$, \inserted{satisfaisant à (72)\textsuperscript{*}} \uline{équivaut} à celle d'un sous-groupe $T_H\subset H$, normalisant $U_0$, tel qu'on ait $T_H\cap U_0=\mathfrak z_H$, $[$que $T_H/\mathfrak z_H$ soit commutatif$]$, \deleted{et} que $\sigma$ normalise $T_H$ et \deleted{\unsure{ind|tant}} satisfasse $\sigma(t)\equiv t^{-1}\ \deleted{\unsure{\mathrm{mod}\,\mathfrak z}}\ \inserted{\mathrm{mod}\,\mathfrak z_H}$, pour $t\in T_H$, et tel enfin qu'on ait $\underbrace{TU_0}_{B_{0H}}\cap U_1=\mathfrak z_H$ (d'où \deleted{\unsure{par}} \inserted{en transformant par $\sigma$} $TU_{\unsure{1}}\cap U_0=\mathfrak z_H$)

% slanted
\margin{* i.e. avec opération \uline{fidèle} de $Z$, sur l'ens $\pi$ des structures de graphe engendré par $(S,\overrightarrow{A_0})$ sous $Z$}

% slanted
\margin{$[\ ]$ résulte du reste}

% slanted
\margin{\uline{NB} On voit alors que \deleted{\unsure{(89)}}}

En effet, si $T_H$ est \deleted{un} sous groupe de $H$ normalisant $U_0$, alors $\inserted{Z\overset{\mathrm{def}}{=}}T_H/U_0\cap T_H\simeq\underbrace{T_H\mathbin{.}U_0}_{B_{0H}}/U_0$, opère \inserted{fidèlement} sur l'espace homogène $S\simeq\deleted{\unsure{G|O}}H/U_0$, il faut vérifier que c'est un groupe d'homothéties, on a \inserted{supposé} la commutativité, d'autre part pour $t\in T_{N}$ on a $\sigma(t)\mathbin{.}t\in\mathfrak z_H\subset U_{1}\cap U_0\inserted{\subset U_1\mathbin{.}U_0}$, un calcul déjà fait, et immédiat nous montre que \unsure{elle|celle|cela} implique que $t$ est une homothétie. \deleted{A}

\deleted{\unsure{Le cas vu|le cas vu} \unsure{alors}, on a utilisé seulement l'inclusion $T_H\cap U_0\subset\mathfrak z$,}

% archive sheet number: 247
\scanpage{249}
\mspage{127}

La condition $TU_0\cap U_1$ ne fait alors qu'exprimer $(72)^*$.

La situation est particulièrement jolie dans le cas où $\unsure{\mathfrak z|\zeta}_H=1$ :

(89) La donnée, pour un graphe homogène avec groupe \uline{simplement} \uline{transitif} $H$ sur $\overrightarrow{A_0}$, défini par $(H,U_0\subset H,\sigma\in H)$, d'un groupe d'homothéties $Z$ satisfaisant la condition de fidélité (72), équivaut à celle d'un sous-groupe $T_H$ de $H$ satisfaisant les conditions

\[
\left\{\begin{aligned}
&\deleted{T_H\text{ commutatif, et norm}}\\
&\deleted{T_H\cap U_0=\{1\}}\\
&\text{a) }\sigma\text{ normalise }T_H\text{ et y induit }t\longmapsto t^{-1}\\
&\text{b) }\deleted{\text{pour}}\ T_H\text{ normalise }U_0\quad(\text{donc \inserted{aussi} }U_1\deleted{\unsure{g|?}}= \sigma(U_0))\\
&\text{c) l'application}\\
&\qquad U_0\times T_H\times U_1\longrightarrow H\\
&\qquad(u_0,t,u_1)\longmapsto u_0tu_1\\
&\text{est \uline{injective}.}
\end{aligned}\right.
\]

\margin{\deleted{$\left\{\begin{aligned}\unsure{T_H\cap U_1}\\[-2pt]\unsure{\ni}U_0\\[-2pt]\unsure{T_H\cap U_1}\end{aligned}\right.$}}

Je vais maintenant expliciter, pour une telle donnée, la construction du diagramme correspondant (64) --- je me borne au cas $H$ simplement transitif sur $\overrightarrow{A_0}$, pour simplifier, cas tellement plus \unsure{joli|jolie} que le cas général !

J'introduis \unsure{le|les} sous-\unsure{groupe|groupes}

\[
\deleted{(90)\qquad U_{0G}=G_{s_0}=\{g\in G\mid gs_0=s_0\}}
\]

\[
\begin{aligned}
\deleted{\unsure{T'_G=}}\ T_G&=G_{(D_0,D_1)}=\{g\in G\mid\exists\lambda,\mu\in Z,\ gs_0=\lambda s_0,gt_0=\mu t_0\}\\
&\qquad\deleted{\text{avec }gD_0=D_0,gD_1=}
\end{aligned}
\tag*{(90)}
\]

Les éléments $\lambda,\mu\inserted{\in Z}$ sont bien déterminés par $g$, et on trouve un hom can. $T_G\longrightarrow Z\times Z$, \deleted{\unsure{évi-}}

% archive sheet number: 248
